\section{Non-Atomic Group Soundness}
\label{sec:group-soundness}

M3-B permits non-atomic blocks and therefore permits partial-block deletions.
The replacement for syntactic block alignment is a semantic reporting
condition: every feasible implementation deletion must remain feasible in the
reference contract after touch-any grouping.

\begin{theorem}[M3-B: non-atomic grouped correctness]
\label{thm:m3b}
Assume active blocks are pairwise disjoint.  If
\(\ResidualFaithfulness\) and \(\GroupSoundness\) hold, then for every grouped
deletion \(G\),
\[
  \GroupedRepair(G)
  \quad\Longleftrightarrow\quad
  \ContractRepair(G).
\]
\end{theorem}

The statement is pointwise predicate equivalence, not an unqualified function
equality.  The Lean declaration is \code{m3b_grouped_correctness}.

\paragraph{Proof route.}
A contract repair supplies a feasible block deletion.  Finite descent chooses
a raw minimal subdeletion without assuming deletion monotonicity of
\(\SatPhi\).  Group soundness makes the grouped deletion contract-feasible,
and contract minimality forces equality with the starting contract repair.
Conversely, every grouped raw report contains a contract-minimal repair; the
outer inclusion-minimality step removes nonminimal grouped images.

\begin{corollary}[Two realizations]
If two realizations of the same reference contract are residually faithful and
group-sound, their grouped repair predicates are pointwise equivalent through
the common contract repair predicate.
\end{corollary}

This corollary is implemented as \code{m3b_two_realization}.  M3-B establishes
grouped correctness, not raw transport.

\paragraph{Candidate B qualification.}
The Candidate B use of M3-B relies on complete artifact-level raw-repair
enumeration and grouped-image verification; the Lean theorem proves the
abstract sufficient condition, not the artifact enumeration.
